\section{K3 Teil 2: Spitzen-Tage ausblenden}
\label{app:report_k3_teil_2}
\subsection*{Setup}
\begin{itemize}
\item \textbf{Lauf:} k3\_\allowbreak{}mask\_\allowbreak{}004
\item \textbf{Stichtage:} SF1: 31.07.2014, SF2: 12.02.2015, SF3: 24.05.2015, SF4: 26.09.2015, SF5: 12.01.2016
\item \textbf{Szenarien:} peak an Tag 0, 8, 17 nach dem Stichtag (Orakel, ohne Einordnung)
\item \textbf{Vergleich:} Ereignistage für Prophet ausgeblendet (mask) gegen nicht ausgeblendet (none), 60 Tage Trainingsfenster
\item \textbf{Commit:} 83412b7 (mit lokalen Änderungen)
\end{itemize}

\begin{landscape}
\subsection*{Differenz des Tagesfehlers mask minus none [W], negativ = Ausblenden besser}
\begin{tablang}[\footnotesize\setlength{\tabcolsep}{3pt}]{W{1.36} V{0.80} V{1.70} V{0.68} V{0.68} V{0.68} V{0.68} V{0.68} V{1.14} W{1.59}}{Szenario & Tage & ohne Ausblenden [W] & SF1 & SF2 & SF3 & SF4 & SF5 & Mittel & Stichtage besser}
peak & 300 & 76,8 & 1,0 & \textminus{}1,1 & 3,1 & \textminus{}0,9 & 0,7 & 0,6 & 2/5 \\
\end{tablang}
\end{landscape}

\subsection*{Hinweise}
\begin{itemize}
\item mask minus none: beide gegen dieselbe Reihe mit Spitzen bewertet, ausgeblendet ist nur das Training
\item Tagesfehler wie in Teil 1: Tage mit unvollständig bewerteten Stunden zählen dort nicht mit
\end{itemize}
